\section{Introduction}
\label{sec:intro}

IoT devices in factories, buildings and hospitals produce a steady stream of computing tasks. Each task can run on a small edge server close to the device or on a large cloud server far away. This range of computing places is called the edge-cloud continuum. Choosing where each task runs is task offloading, and it decides whether the task finishes within its latency limit \cite{mao2017}. A service level agreement (SLA) states this limit for each type of task. When a task finishes late, the SLA is violated.

IoT workloads are mixed. A control-loop signal from a machine may need an answer within tens of milliseconds. A batch of sensor readings for analytics can wait longer. Both kinds of tasks compete for the same edge capacity. When the edge is busy, an offloading policy has to decide which tasks it protects.

Deep reinforcement learning (RL) \cite{mnih2015,sutton2018} is a common way to learn offloading policies. The agent observes the system, places a task, and receives a reward that says how good the placement was. Many RL offloading methods give every missed task the same penalty, whatever its class \cite{nieto2024}. With such a reward the average can look good while a small class of critical tasks keeps missing its limit. A natural fix is to give each class its own penalty. Recent SLA-aware schedulers do this, for example with Gold, Silver and Bronze tiers \cite{yamsani2026}. They change the penalty together with the state, the action space and the network, and they compare against rules and a class-blind agent only. Two questions stay open. What does the per-class penalty do on its own? And how far is any of these policies from the best policy that is possible?

This paper answers both questions with an exact benchmark. The offloading systems studied here are small enough that their full state space can be listed, yet they have heterogeneous edge nodes, up to three task classes, random background traffic and a cloud that meets some limits only by chance. Finite-horizon dynamic programming (DP) then gives the optimal policy for any cost, the best possible trade-off between the classes, and the exact expected performance of any policy, including a trained neural network. Every agent and every rule in this paper is placed against the optimum, and no reported number depends on a finite test set.

The paper answers five questions.
Q1: With a perfect learner, how much does a per-class penalty change the result, compared with a class-blind penalty of the same average, and why?
Q2: Do trained PPO, DQN and A2C agents show the same effect, and how far are they from the optimum?
Q3: How do the agents compare with rules that know the latency model and are tuned on the exact model, and with published SLA-aware reward designs?
Q4: Does the effect hold in larger systems with heterogeneous nodes and three classes, where tuned rules are far from the optimum?
Q5: What happens when the system changes after rules and agents were built?

The contributions are the following. First, the paper gives an exact model of a family of multi-class edge-cloud offloading systems with sparse transitions, shows that its observation is a full state, and uses it to compute optimal policies, Pareto frontiers and exact scores of trained agents. The model is checked against the simulator in the open test suite. Second, it proves that a reward with a bonus for a met limit ranks policies by the per-class penalty plus the bonus, not by the penalty alone. This explains why a class-weighted reward of 10 and 2 lands 1\% above the optimum, and gives a corrected penalty rule whose optimal policy is the optimum. It also shows that a class-blind reward minimizes the overall miss rate, so its distance to the optimum grows without bound as the classes differ more. Third, it reports a 10-seed study of fairly tuned PPO, DQN and A2C agents with class-weighted and class-blind rewards, and of PPO with the corrected reward, against rules tuned on the exact model and against two published reward designs, in four systems and under five changes of the system. Fourth, all code, trained agents and results are open.

The rest of the paper is organized as follows. Section~\ref{sec:related} reviews related work and states how each paper differs from this one. Section~\ref{sec:model} describes the systems and the rewards. Section~\ref{sec:exact} describes the exact model. Section~\ref{sec:reward} analyzes the rewards exactly. Section~\ref{sec:method} describes the agents, the baselines and the protocol. Section~\ref{sec:results} reports the results. Section~\ref{sec:discussion} discusses them and the limits of the study, and Section~\ref{sec:conclusion} concludes.

\section{Related Work}
\label{sec:related}

\subsection{Deep RL for Task Offloading}
Mao et al. \cite{mao2017} survey mobile edge computing from the communication side and describe offloading as a joint choice of where a task runs and how its data is sent. Most work since then learns this choice with deep RL. Hu et al. \cite{hu2021} learn task scheduling, transmit power and edge resources together for power IoT, with an objective that is a weighted sum of latency and energy. Chen and Li \cite{chen2025} combine offloading with wireless power transfer, and Zhang et al. \cite{zhang2024} combine it with reconfigurable intelligent surface (RIS) beamforming in a cloud-edge system. These three works decide how fast data reaches the edge and how energy is spent. This paper takes the network side as fixed and asks which class is protected once the edge is full.

\subsection{QoS and SLA-Aware Offloading}
AlQerm and Pan \cite{alqerm2021} maximize a single quality-of-experience score across many applications. Heidari et al. \cite{heidari2023} maximize a single efficiency score that mixes delay, energy and failure rate. Both works fold all tasks into one number. This paper reports a separate SLA violation rate for each class, together with the full trade-off curve between the classes.

Nieto et al. \cite{nieto2024} train actor-critic agents in a 5G edge-cloud system with three task classes in the state. Every class receives the same fixed penalty when a task fails, so the agent sees the class but is not told that some classes matter more. This design is the class-blind reward of this paper. SLA-ORECS \cite{lan2024} guarantees SLAs by moving resources between channels, not by learning a placement for each task. Yamsani and Chenna Reddy \cite{yamsani2026} train a DQN scheduler with an SLA violation risk score in the state, Gold, Silver and Bronze penalty tiers that grow when a tier misses its target, and pruning of unsafe actions. Their ablation removes the whole SLA reward at once, so the effect of the tiers alone is not measured, and their policies are not compared with an optimum. This paper rebuilds their reward, risk features and pruning in its own systems and scores them exactly or with large simulations.

\subsection{Reward Design and Constrained RL}
Ng et al. \cite{ng1999} show which reward changes keep the optimal policy unchanged. A constant bonus for a met limit is not such a change when penalties differ by class, and Section~\ref{sec:reward} shows exactly how it moves the optimal policy. Constrained Markov decision processes \cite{altman1999} and Lagrangian RL methods \cite{tessler2019} raise a penalty while a constraint is violated. The adaptive tiers of \cite{yamsani2026} follow the same idea with a fixed step size.

\subsection{Optimal Control of Multi-Class Queues}
For a single pool of servers and several classes, the optimal admission policy reserves capacity for the more valuable class, which gives a threshold or reservation rule \cite{miller1969}. The tuned reservation rules of this paper follow this result. With heterogeneous servers, random background traffic and a cloud that sometimes meets a limit, no simple optimal form is known, and this paper computes the optimum numerically.

\subsection{Optimal Benchmarks and Evaluation of Deep RL}
For small Markov decision processes, DP gives the optimal policy and the exact value of any policy \cite{puterman1994,sutton2018}. Offloading studies rarely use such a benchmark, because realistic systems are too large to solve. This paper keeps the systems small enough to solve, so that every learned result can be read against the optimum. Deep RL results also vary strongly between random seeds and with the evaluation protocol \cite{henderson2018,agarwal2021}. This paper uses 10 seeds, reports confidence intervals and effect sizes, corrects for multiple tests, tunes every algorithm with the same budget, and removes evaluation noise with the exact model.
